\phantomsection\bheading{section}{7.\quad Interaction with the rest of beni}{section}{\protect\numberline{7.}Interaction with the rest of beni}

\phantomsection\bheading{subsection}{7.1\quad The checker: a pass after type checking}{subsection}{\protect\numberline{7.1}The checker: a pass after type checking}

The ownership pass runs where \code{Effects.run} runs — after P5, with every type solved, before P6 and publication (\code{checker-v2.md} §26) — and for the same reason Futhark gave when it moved alias checking out of its type checker in 2026: the old design “had always been hindered by having to work with incomplete type information”, and a separate pass over “a fully well-typed program … simply had to figure out whether any of the in-place constraints were violated (essentially, whether in-place updates are semantically observable)” (2026 rewrite post, ll. 110–117). Its messages carry the solved types, and it never touches \code{unify} or the store. It needs one thing the checker does not keep today: the solved type of every BIR instruction, or at least of every list-typed one, to decide paths and crossing. The write-set pass already asks the checker for “paths [that] are typed, which the checker's solved types make possible” (\code{write-sets.md} §2.1), so the hook exists.

It also needs a per-instruction \emph{cell site} id (the instruction index, as the write-set pass uses the instruction for its index symbols) and, for messages, the source span of every holder, which the instruction's \code{main\_\allowbreak{}token} gives (\code{frontend.md} §3.6).

\phantomsection\bheading{subsection}{7.2\quad Effect bits}{subsection}{\protect\numberline{7.2}Effect bits}

The two inferred bits and the ownership classes ride on the same function-type sites, are solved in the same per-module phase and are published in the same kind of block. They interact in one place: a function that \textbf{suspends} parks its locals in a continuation (\code{transparent-effects- proposal.\allowbreak{}md} §16); the continuation is one-shot (A7), so a parked local is the same holder it was, and liveness is unchanged. A \code{sync} demand says nothing about ownership. An \code{impure} call is a call like any other for holders; the optimiser's duty to keep every impure call in place (\code{language.md} §6) is what A5 relies on.

\phantomsection\bheading{subsection}{7.3\quad Static dispatch}{subsection}{\protect\numberline{7.3}Static dispatch}

A method call's callee is the dispatch table's term (\code{checker-v2.md} §13.1); its summary is the callee's. A \code{param} term (evidence the enclosing function receives) is a class instantiated at the caller, §3.7. Derived \code{eq}/\code{compare} are read-only ($\mathsf{borrowed}$). The \code{where} suffix and the summary block are both inferred facts in the interface, and both churn the same way; research 19 §4's instrument can measure the second as it measured the first.

\phantomsection\bheading{subsection}{7.4\quad Write sets}{subsection}{\protect\numberline{7.4}Write sets}

Two reads and one restatement:

\begin{itemize}
\item \textbf{The checker computes the entry state itself} (P5) and reads from the write-set pass only the program's shape: which declaration is \code{update}, which \code{init}, the key tree. The escape scan \code{browser-direct} §11 item 4 lists for S7 is this computation, done once here for both; the write-set interpreter does not analyse commands (\code{write-sets.md} §3.3, \emph{Effects}) and records no aliasing between two model paths, so it cannot supply it.
\item \textbf{The write-set consumer reads nothing new}, but its soundness statement must be read with one word changed. \code{write-sets.md} §1.3 defines “changed” as “not the same value: not \code{===}”. Under in-place writes a \code{value} write at $\rho .\mathit{rows}[k]$ leaves $M'@\rho .\mathit{rows}\mathrel{\mathtt{===}}M@\rho .\mathit{rows}$. The write set still \emph{names} the path (the \code{Lst(p, kept, [k], …)} row is read off the arm's text, not off identity), so a consumer that calls the groups the write set names is still correct; what it may no longer do is conclude “unchanged” from \code{===} on a list path — which is exactly \code{browser-direct} §7.4's first rule and §6.2's tag guards. The restatement: \emph{“changed” means “may differ in content from what was there before the dispatch”}, and \code{write-sets.md} §5.4's requirement 1 (“an in-place list write never touches a list a program still holds”) becomes this checker's theorem rather than a requirement on the emitter.
\item \textbf{The write-set interpreter is the right engine for the model paths, not for the general checker.} It is whole-program, per message key, and tracks identity; the ownership checker is per function, directional over cells, and publishes summaries. They meet at \code{update}'s entry (P5) and nowhere else.
\end{itemize}

\phantomsection\bheading{subsection}{7.5\quad The backend: in-place emission, and what replaces the trie}{subsection}{\protect\numberline{7.5}The backend: in-place emission, and what replaces the trie}

Under the rule every accepted demand is in place, so the backend needs no per-site decision: it emits each writer's call as it does today, and \code{core/List}'s writers \emph{are} the in-place operations:

{\small\setlength{\emergencystretch}{3em}
\begin{longtable}{>{\raggedright\arraybackslash}p{\dimexpr0.250\linewidth-2\tabcolsep\relax}>{\raggedright\arraybackslash}p{\dimexpr0.300\linewidth-2\tabcolsep\relax}>{\raggedright\arraybackslash}p{\dimexpr0.450\linewidth-2\tabcolsep\relax}}
\toprule
\textbf{primitive} & \textbf{today (\code{backend.md} §4)} & \textbf{under the rule} \\
\midrule
\endfirsthead
\toprule
\textbf{primitive} & \textbf{today (\code{backend.md} §4)} & \textbf{under the rule} \\
\midrule
\endhead
\bottomrule
\endfoot
\code{set xs i v}, \code{update}, \code{swap} & copy $\le$ 256, else trie path copy & \code{a[i] = v} on the base; O(1) \\[0.35em]
\code{push xs v} & copy < 32, else tail claim & \code{a.push(v)}; amortised O(1) \\[0.35em]
\code{pop xs} & copy $\le$ 256, else header & \code{a.pop()}; O(1) \\[0.35em]
\code{insertAt}, \code{removeAt} & O(n) plain copy & \code{a.splice(i, …)}; O(n) memmove, what vanilla writes \\[0.35em]
\code{cons x xs} (\code{[ x, …xs ]}) & head claim / $\le$ 31 copy / convert & \code{a.unshift(x)}; \textbf{O(n)} (§8) \\[0.35em]
\code{append xs ys} & push onto a trie or concatenate & \code{for (y of ys) a.\allowbreak{}push(y)}; $O(\lvert\mathit{ys}\rvert)$ \\[0.35em]
\code{copy xs} & — & \code{a.slice()} \\[0.35em]
a write through a view $(b,\allowbreak{}o)$ & — & the same on $b$ at $o+i$; a new header for \code{push}/\code{pop} \\[0.35em]
\bottomrule
\end{longtable}}

The \textbf{trie goes}: E1tp's claims, the radix offset, the head and tail buffers, \code{\$plain()}'s cache and the thresholds (\code{backend.md} §4, \emph{E1t}, \emph{E1tp}) exist to make writes to \emph{shared} lists cheap, and under the rule no shared list is ever written. What stays: the plain form, the view form (for a pattern's \code{rest}, with the scalar-view loop rewrite still removing most of them), the builder (\code{backend.md} §8), and the three-point protocol for readers (\code{length}, \code{Array.isArray}, \code{\$plain()}, the last now only for views). Research 42 §7.4 measured the sibling at 706 bytes brotli without the trie against 1 514 with it, and research 40 measured the write half at \~{}450 B in context; \code{browser-direct} §7.2 found the table app shipping the write half for a list it only reads. All of that is removed. \code{List.beni}'s loops are unchanged.

What the backend must still keep: the identity promises it chooses to keep (§10, change 2), the builder's linearity (invariant 5), and the fact that \code{Basics.append} on \code{appendable} stays a sibling copy (\code{language.md} §6.8) — \code{++} on lists is \code{List.append}, a demand on its first argument, which §10 (change 4) asks the owner to confirm.

\phantomsection\bheading{subsection}{7.6\quad The release optimiser}{subsection}{\protect\numberline{7.6}The release optimiser}

The verdict is a function of the source; \code{--release} emits the same in-place calls as the development build, so “a release build behaves exactly as the development build does” holds with no new exception — \textbf{provided the optimiser treats a demand as a barrier} (A5). Today's item 1 (\code{backend.md} §9) folds a single-use pure binding into its use across other such bindings, because \code{language.md} §6's 2026-10-02 amendment lets evaluations that make no call commute; a read of a cell and a demand on it do not commute, so the demand summaries become an input to \code{Opt} as the \code{impure} bit is, and a read of a cell is never moved across a demand on it. With that, the rest of the contract carries over: an unused binding whose right-hand side is a pure call may be dropped (an unused \code{List.set} is then not performed, which no reader can tell), two surviving calls are never reordered, inlining substitutes bodies and never duplicates an argument's evaluation. Single-use inlining may inline a $\mathit{once}$ closure's body into its one call. The specialiser's copies inherit the general verdict (§5.16). \code{Minify} and \code{Rename} touch names, not writes. \code{derived\_\allowbreak{}bytes} and reachability are untouched: \code{copy} ships only when reached.

\phantomsection\bheading{subsection}{7.7\quad Records in place (\code{browser-direct} S7)}{subsection}{\protect\numberline{7.7}Records in place (\code{browser-direct} S7)}

Not in this thesis's scope, but the same machinery applies: a record path is a “cell” whose demand is the record-update-in-place the lowering would emit; its holders are the same places; \code{language.md} §11.12's identity promise becomes the only difference (a written record keeps its identity, which §7.4's rule 1 already forbids the renderer from relying on). The one hazard is that records \emph{are} compared by identity by the renderer everywhere (the per-hole reference check, rule 8), so record paths need a stronger A4 than lists do; §7.4 (there) already states the two rules.

\phantomsection\bheading{subsection}{7.8\quad Incremental builds}{subsection}{\protect\numberline{7.8}Incremental builds}

The summary block is in the interface and the hash; the firewall fires when it changes. Two consequences, both to be measured (§9): churn (a body edit that turns a $\mathsf{borrowed}$ parameter $\mathsf{consumed}$ re-checks every importer — necessary, since their K-checks change), and the warm budget (\code{fast-compiler.md} §2: 15 ms for a body edit, 60 ms for a signature edit). A body edit that changes a summary \emph{is} a signature edit under this design, and \code{plans/queue.md} row 55 records that signature edits do not yet meet their budget for other reasons.

\phantomsection\bheading{subsection}{7.9\quad Determinism}{subsection}{\protect\numberline{7.9}Determinism}

§3.10. One addition for the dump: cells print as \code{site <decl>:<inst>} and \code{param <i>.<path>}, both source-derived; summaries print in scheme-site order; holders in source order. The \code{--jobs=1}/\code{--jobs=8} scenario byte-compares the ownership dump beside the others.
